\section*{Chapitre \ref{ch:b3:virology} --- Virologie}

\begin{solution}{exo:b3:virology:1}
Polio~: IV, ARN à brin positif --- une ARN polymérase ARN-dépendante.
Grippe~: V, ARN à brin négatif --- sa propre ARN polymérase, emportée
dans le \omterm{def:b3:virology:virus}{virion}. VIH~: VI --- la transcriptase inverse (et
l'intégrase). Herpès simplex~: I, ADN double brin --- l'hôte pourrait
en principe tout fournir, mais il apporte sa propre ADN polymérase et
sa thymidine kinase. Hépatite B~: VII --- une transcriptase inverse
pour copier son ARN prégénomique en ADN. Rotavirus~: III, ARN double
brin --- une ARN polymérase ARN-dépendante dans la particule.
\end{solution}

\begin{solution}{exo:b3:virology:2}
L'ADN du phage ($^{32}$P) entrait dans les bactéries et la protéine
($^{35}$S) restait dehors et pouvait être arrachée~; la descendance du
phage héritait du $^{32}$P. C'est donc l'ADN que le phage injecte et
qui dirige la fabrication de nouveaux phages. Si la protéine avait
porté les instructions, on aurait trouvé le $^{35}$S à l'intérieur et
transmis à la descendance.
\end{solution}

\begin{solution}{exo:b3:virology:3}
$84/(0.1\times 10^{-7}) = 8.4\times 10^{9}$ unités formant plage par
mL. Le microscope compte chaque particule, y compris les \omterm{def:b3:virology:virus}{capsides}
vides, les particules à génome défectueux et celles qui n'arrivent ni à
s'attacher ni à injecter~; seules les particules qui mènent une
infection à son terme font des plages.
\end{solution}

\begin{solution}{exo:b3:virology:4}
Attachement~: inhibiteurs d'entrée (maraviroc), anticorps
neutralisants. Entrée et fusion~: inhibiteurs de fusion, bloqueurs de
l'acidification endosomale. Décapsidation~: médicaments qui se fixent à
la \omterm{def:b3:virology:virus}{capside} (pléconaril~; lénacapavir). Expression et réplication~:
\omterm{def:b3:virology:drugs}{analogues nucléosidiques}, inhibiteurs de la transcriptase inverse et de
l'intégrase, interféron, \omterm{def:b3:rna-regulation:rnai}{interférence par ARN}. Assemblage~: \omterm{def:b3:virology:drugs}{inhibiteurs
de protéase} (des polyprotéines non clivées ne peuvent pas s'assembler).
Libération~: inhibiteurs de neuraminidase (grippe), et lymphocytes~T
cytotoxiques qui tuent la cellule avant la libération.
\end{solution}

\begin{solution}{exo:b3:virology:5}
Une protéine de \qty{30}{kDa} fait $30\,000/110 = 273$ résidus, soit
$818$ nucléotides~: \qty{0.82}{kb}, \qty{18}{\%} d'un génome de
\qty{4.5}{kb}. Coder séparément les $180$ copies demanderait
\qty{147}{kb}, trente-trois fois le génome entier~; la protéine de
\omterm{def:b3:virology:virus}{capside} elle-même (\qty{5.4}{MDa}) pèse quatre fois le génome
(\qty{1.5}{MDa}).
\end{solution}

\begin{solution}{exo:b3:virology:6}
$V = 10^{5}(3e^{-0.5t} - 0.5e^{-3t})/2.5$~: $t = 0.5$~: $8.9\times
10^{4}$~; $t = 2$~: $4.4\times 10^{4}$~; $t = 10$~: $8\times 10^{2}$.
Sous $50$ quand $1.2\times 10^{5}e^{-0.5t} = 50$~: $t = 2\ln(2400) \approx
\qty{16}{d}$.
\end{solution}

\begin{solution}{exo:b3:virology:7}
$uL = 0.36$~; fraction sans erreur $e^{-0.36} = 0.70$. Trois fois plus~:
$uL =
1.08$, $e^{-1.08} = 0.34$ --- la plupart des descendants portent
désormais des mutations et la séquence la mieux adaptée est diluée vers
le seuil. Un coronavirus de \qty{30}{kb} avec $u = 10^{-4}$ aurait
$uL = 3$ et seulement \qty{5}{\%} de copies sans erreur, au-delà du
seuil~; son exonucléase ramène $u$ à $10^{-5}$, $uL = 0.3$, trois
quarts sans erreur.
\end{solution}

\begin{solution}{exo:b3:virology:8}
Glissement~: les mutations ponctuelles de l'hémagglutinine et de la
neuraminidase s'accumulent sous la sélection des anticorps, de sorte
que la souche de chaque saison est en partie nouvelle et que l'immunité
de l'an dernier est en partie périmée. Cassure~: la co-infection d'une
même cellule par une souche humaine et une souche animale réassortit
les huit segments, et un sous-type d'hémagglutinine venu des oiseaux ou
des porcs apparaît dans un \omterm{def:b3:virology:virus}{virus} adapté à l'homme. Personne n'a
d'anticorps contre le nouveau sous-type, de sorte que toute la
population est sensible d'un coup --- la condition d'une pandémie,
remplie en 1918, 1957, 1968 et 2009 par une hémagglutinine nouvelle à
chaque fois.
\end{solution}

\begin{solution}{exo:b3:virology:9}
La thymidine kinase de l'herpèsvirus phosphoryle l'aciclovir, que les
kinases cellulaires touchent à peine~; les kinases cellulaires achèvent
ensuite le triphosphate, que l'ADN polymérase virale incorpore et qui,
faute d'hydroxyle 3$'$, termine la chaîne. Les cellules non infectées
ne fabriquent jamais la forme active, et la polymérase virale la
préfère à la cellulaire. Résistance~: perte ou altération de la
thymidine kinase virale (la plus fréquente), ou mutations de la
polymérase virale qui excluent l'analogue.
\end{solution}

\begin{solution}{exo:b3:virology:10}
À $T$ constant les équations sont linéaires en $(I, V)$, de matrice
$\begin{pmatrix} -\delta & \beta T \\ p & -c \end{pmatrix}$, de trace
$-(\delta + c) < 0$ et de déterminant $\delta c - \beta Tp$. La trace
des valeurs propres étant négative, l'une est positive exactement quand
le déterminant est négatif~: $\beta Tp > c\delta$, c'est-à-dire
$R_{0} = \beta
Tp/(c\delta) > 1$. Les facteurs~: $p/\delta$ est le nombre de \omterm{def:b3:virology:virus}{virions}
qu'une cellule infectée fabrique dans sa vie~; $\beta T/c$ est le
nombre de cellules qu'un \omterm{def:b3:virology:virus}{virion} infecte dans la sienne~; leur produit
est le nombre de cellules infectées qu'une cellule infectée produit.
\end{solution}

\begin{solution}{exo:b3:virology:11}
$10^{6} = 2^{20}$~: vingt demi-vies, $880$ mois, environ $73$ ans ---
plus qu'une vie. Les médicaments qui bloquent la réplication ne
touchent pas un provirus qui n'est pas transcrit~; le réservoir
persiste, et à l'arrêt du traitement une seule cellule qui se réactive
relance l'infection. Une guérison doit retirer ou faire taire
définitivement le réservoir~: tuer les cellules latentes (les réactiver
sous traitement pour qu'elles meurent ou soient tuées), exciser ou
désactiver le provirus, ou remplacer le système immunitaire par un
système que le \omterm{def:b3:virology:virus}{virus} ne peut pas infecter.
\end{solution}

\begin{solution}{exo:b3:virology:12}
$uL = 0.03$ mutation par génome~; $10^{9}$ génomes par jour en portent
$3\times
10^{7}$ parmi $9\times 10^{4}$ mutations ponctuelles possibles~: chacune
apparaît environ $300$ fois par jour. Un double mutant précis apparaît
à $10^{-12}$ par génome, soit $10^{-3}$ par jour --- une fois en trois
ans~; un triple à $10^{-18}$, jamais. Deux médicaments à mutations de
résistance indépendantes suffiraient en principe, trois donnent une
marge en cas d'observance imparfaite. Deux médicaments qui se fixent à
la même enzyme peuvent être défaits par une seule mutation qui modifie
le site ou la conformation de l'enzyme, et leurs mutations de
résistance ne sont pas indépendantes~; ils comptent pour moins de deux.
\end{solution}

\begin{solution}{pb:b3:virology:1}
\textbf{1.} $2\times 10^{5}\times 1.5\times 10^{4} = 3\times 10^{9}$
\omterm{def:b3:virology:virus}{virions}.
\textbf{2.} Éliminés $cV = 3\times 3\times 10^{9} \approx 10^{10}$ par
jour~; produits autant.
\textbf{3.} $I^{*} = cV^{*}/p = 9\times 10^{9}/500 = 1.8\times 10^{7}$
cellules infectées~; il en meurt $\delta I^{*} = 9\times 10^{6}$ par
jour.
\textbf{4.} ARN $2\times 9700\times 330 = 6.4\times 10^{6}$ Da~;
protéine $2000\times 25\,000 = 5\times 10^{7}$ Da~; total
$5.6\times 10^{7}$ Da $\approx 10^{-16}$ g. Tous les \omterm{def:b3:virology:virus}{virions} ensemble~:
$3\times 10^{9}\times
10^{-16} = 3\times 10^{-7}$ g, un tiers de microgramme --- trois mille
fois moins qu'un grain de sable.
\textbf{5.} $1.8\times 10^{7}/10^{12} = 2\times 10^{-5}$ du stock. Dix
ans à $9\times 10^{6}$ morts par jour font $3\times 10^{10}$ cellules,
soit trois pour cent du stock~: le stock cède non par une soustraction
arithmétique, mais parce qu'une décennie de régénération forcée et
d'activation chronique épuise la capacité à les remplacer.
\textbf{6.} $9\times 10^{6}/24 \approx 4\times 10^{5}$ cadavres en cours
d'élimination à tout instant.
\textbf{7.} $V = 2\times 10^{5}(3e^{-0.5t} - 0.5e^{-3t})/2.5$~: $t = 1$~:
$1.4\times 10^{5}$~; $t = 3$~: $5.4\times 10^{4}$~; $t = 7$~: $7\times
10^{3}$~; $t = 14$~: $2\times 10^{2}$.
\textbf{8.} $2.4\times 10^{5}e^{-0.5t} = 50$~: $t = 2\ln(4800) \approx 17$
jours.
\textbf{9.} Les jours 3 à 10 sont sur la phase lente~: la pente de
$\ln V$ y vaut $-\delta$ (de $5.4\times 10^{4}$ à $1.6\times 10^{3}$ en
sept jours, $\delta = 0.50$). La phase rapide ne dure que quelques
heures ($1/c$), de sorte que $c$ exige des prélèvements toutes les
quelques heures le premier jour, ajustés à la formule à deux
exponentielles.
\textbf{10.} Les cellules infectées de façon latente, qui libèrent du
\omterm{def:b3:virology:virus}{virus} en se réactivant~; leur vitesse de décroissance vaut $\ln 2/44$
par mois, soit environ $5\times 10^{-4}$ par jour.
\textbf{11.} \omterm{def:b3:virology:virus}{Virion} $\ln 2/3 = \qty{0.23}{d}$, soit environ
\qty{5.5}{h}~; cellule infectée $\ln 2/0.5 = \qty{1.4}{d}$.
\textbf{12.} Une charge constante avait l'air d'un \omterm{def:b3:virology:virus}{virus} qui ne faisait
rien~; les équations y montrent un état stationnaire où $10^{10}$
\omterm{def:b3:virology:virus}{virions} sont fabriqués et éliminés et $10^{7}$ lymphocytes~T infectés
et tués chaque jour.
\textbf{13.} $9700\times 10^{-5} \approx 0.1$ mutation par génome~;
$10^{10}\times 0.1 = 10^{9}$ mutations par jour.
\textbf{14.} $3\times 9700 \approx \num{29000}$ mutations ponctuelles
possibles~; chacune apparaît $10^{9}/\num{29000} \approx 3\times 10^{4}$
fois par jour.
\textbf{15.} Double précis~: $10^{-10}$ par génome, environ un par jour
sans traitement~; triple précis~: $10^{-15}$, $10^{-5}$ par jour ---
une fois en trois siècles.
\textbf{16.} $10^{7}\times 10^{-15} = 10^{-8}$ par jour, $4\times 10^{-6}$
par an.
\textbf{17.} Un médicament disparu, des mutants résistants aux deux
autres suffisent~: un double précis à $10^{-10}$~; $7\times 10^{7}$
génomes dans la semaine en donnent $7\times 10^{-3}$ attendus ---
encore improbable, sauf si la charge remonte vers $10^{10}$ par jour
pendant le manquement, auquel cas cela devient un par jour et
l'échappement est probable. Les prises manquées sont la porte d'entrée
de la résistance.
\textbf{18.} Les analogues partagent le site actif de la transcriptase
inverse, et une seule mutation de ce site (ou un jeu de mutations aux
analogues de la thymidine) change la façon dont l'enzyme traite
plusieurs d'entre eux~: la résistance croisée. Deux nucléosides
comptent pour moins de deux médicaments indépendants, de sorte que les
protocoles associent des classes --- un nucléoside, un inhibiteur
d'intégrase, un \omterm{def:b3:virology:drugs}{inhibiteur de protéase} ou un non-nucléoside.
\textbf{19.} $20$ demi-vies, $880$ mois, $73$ ans~; jusqu'à $10^{4}$,
$\log_{2}100 = 6.6$ demi-vies, $290$ mois, $24$ ans.
\textbf{20.} D'un \omterm{def:b3:virology:virus}{virion} à $3\times 10^{9}$ à $R_{0} = 8$ par
génération~: $\ln(3\times 10^{9})/\ln 8 = 10.5$ générations, trois
semaines.
\textbf{21.} Réveiller et tuer~: réactiver les cellules latentes sous
traitement pour qu'elles meurent ou soient tuées --- obstacle~: aucun
agent ne les réactive toutes, et les cellules réactivées ne sont pas
tuées de façon fiable. Édition génomique~: exciser le provirus ou
détruire le récepteur CCR5 dans les cellules du patient --- obstacle~:
atteindre chaque cellule. (Remplacer la moelle par des cellules de
donneur dépourvues de CCR5 a guéri une poignée de patients, à un risque
acceptable seulement lorsqu'une greffe était de toute façon
nécessaire.)
\textbf{22.} $1 - 1/R_{0}$~: \qty{67}{\%} pour $R_{0} = 3$,
\qty{80}{\%} pour $5$.
\textbf{23.} Le traitement met $\beta \to 0$ chez le patient, la charge
tombe par le théorème à presque rien, et la transmission, qui est
proportionnelle à la charge, cesse~; traiter toute personne infectée
fait passer le $R_{0}$ de la population sous un.
\textbf{24.} $f = (1 - 1/R_{0})/E$~: avec $R_{0} = 4$ et $E = 0.7$, $f =
0.75/0.7 = 1.07$ --- plus que tout le monde~: le \omterm{def:b3:virology:drugs}{vaccin} seul ne peut
pas atteindre le seuil~; avec $E = 0.9$, $f = 0.83$.
\textbf{25.} Environ $10^{10}$ \omterm{def:b3:virology:virus}{virions} produits par jour~; indétectable
vers le jour $17$~; quelque $4\times 10^{-6}$ triple mutant par an sous
traitement.
\end{solution}
